\chapter{Ciało, pierścień, przestrzeń mierzalna}
\label{chap:cialo-pierscien-przestrzen-mierzalna}
% Krótszy nagłówek bieżący mieści pełną nazwę rozdziału w szerokości strony.
\markright{\MakeUppercase{\thechapter. Ciało, pierścień, przestrzeń mierzalna}}

\section{Ciało i sigma-ciało zbiorów}
\label{sec:cialo-sigma-cialo}

\begin{definicja}[Ciało zbiorów]
\label{def:cialo}
Ciałem zbiorów w przestrzeni $X$ nazywamy rodzinę $\A\subseteq\mathcal{P}(X)$
spełniającą następujące warunki:
\begin{enumerate}
  \item $\varnothing\in\A$;
  \item jeżeli $A,B\in\A$, to $A\cup B\in\A$;
  \item jeżeli $A\in\A$, to $A^c\in\A$, gdzie $A^c=X\setminus A$.
\end{enumerate}
\end{definicja}

\begin{definicja}[Sigma-ciało zbiorów]
\label{def:sigma-cialo}
Sigma-ciałem (\mbox{$\sigma$-ciałem}) zbiorów w przestrzeni $X$ nazywamy
rodzinę $\A\subseteq\mathcal{P}(X)$, dla której warunki 1) i 3) z definicji
ciała są zachowane, a warunek 2) zastępujemy warunkiem:
\begin{enumerate}[label=2\textquotesingle),start=2]
  \item jeżeli $A_n\in\A$ dla każdego $n\geq1$, to
  \[
    \bigcup_{n=1}^{\infty}A_n\in\A.
  \]
\end{enumerate}
\end{definicja}
